\chapter{Solvability, nilpotence, and semisimplicity}
\label{ch:structure}

\section{Common invariant vectors}

Throughout this chapter the Lie algebras and the spaces acted on
are finite dimensional. Engel's theorem is valid over any field.
Lie's theorem and the structure theorems that follow it use
characteristic zero.

\begin{theorem}[Engel]\label{res:B01}
Let \(0\ne V\) and let \(\mathfrak a\subseteq\mathfrak{gl}(V)\)
be a Lie subalgebra every element of which is a nilpotent
endomorphism. There exists \(0\ne v\in V\) with \(\mathfrak a v=0\).
There is a basis in which all elements of \(\mathfrak a\) are
strictly upper triangular.

A finite-dimensional Lie algebra \(\mathfrak g\) is nilpotent
if and only if \(\operatorname{ad}x\) is nilpotent for every
\(x\in\mathfrak g\).
\end{theorem}

\begin{proof}
Induct on \(\dim\mathfrak a\) to prove the common-kernel
assertion for every representation by nilpotent operators.
The zero algebra is immediate. Choose a maximal proper
subalgebra \(\mathfrak b\) of a nonzero \(\mathfrak a\).
For \(b\in\mathfrak b\), nilpotence of \(b\) on \(V\)
implies nilpotence of \(\operatorname{ad}b\) on
\(\operatorname{End}(V)\): if \(b^N=0\), expand
\((L_b-R_b)^{2N-1}\) and use commutativity of \(L_b,R_b\).
Its induced action on \(\mathfrak a/\mathfrak b\) is
therefore nilpotent. The acting image has dimension at most
\(\dim\mathfrak b<\dim\mathfrak a\), so induction gives
an \(x\notin\mathfrak b\) with
\([\mathfrak b,x]\subseteq\mathfrak b\).
The subspace \(\mathfrak b+kx\) is a subalgebra.
Maximality implies it is all of \(\mathfrak a\); hence
\(\mathfrak b\) is an ideal of codimension one.

Induction on \(\mathfrak b\) acting on \(V\) gives
\(W=\{v:\mathfrak b v=0\}\ne0\).
For \(v\in W\) and \(b\in\mathfrak b\),
\(b(xv)=x(bv)+[b,x]v=0\), so \(W\) is \(x\)-stable.
The nilpotent operator \(x|_W\) has a nonzero kernel.
A vector in that kernel is killed by all of \(\mathfrak a\).
Applying the assertion on successive quotients by the
resulting one-dimensional invariant subspaces gives a full
flag on which every diagonal action is zero. This is the
strict triangularization.

If \(\gamma_{c+1}\mathfrak g=0\), each
\((\operatorname{ad}x)^c\) is zero. Conversely, apply strict
triangularization to the adjoint image on \(\mathfrak g\).
The product of \(\dim\mathfrak g\) strictly upper triangular
operators is zero. Such products span the iterated brackets
defining the lower central series, so that series terminates.
\end{proof}

\begin{lemma}[The cyclic-space calculation]\label{supp:cyclic-weight}
Assume characteristic zero. Let \(\mathfrak b\) be an ideal
of \(\mathfrak b+kx\), and suppose \(0\ne v\in V\) satisfies
\(bv=\lambda(b)v\) for all \(b\in\mathfrak b\).
The space \(W=k[x]v\) is stable under both \(x\) and
\(\mathfrak b\), and each \(b\) has the constant diagonal
\(\lambda(b)\) on \(W\). Moreover,
\(\lambda([x,\mathfrak b])=0\), and the entire common
\(\lambda\)-eigenspace of \(\mathfrak b\) is \(x\)-stable.
\end{lemma}

\begin{proof}
For every \(b\), induction on \(j\), using
\(bx^jv=x(bx^{j-1}v)+[b,x]x^{j-1}v\), gives
\[
 bx^jv-\lambda(b)x^jv
       \in\operatorname{span}(v,xv,\ldots,x^{j-1}v).
\]
The induction is simultaneous for all \(b\), since
\([b,x]\in\mathfrak b\).
Let \(m\) be the first integer for which \(x^mv\)
depends on earlier powers. The first \(m\) powers form
a basis of \(W\), proving stability and the diagonal claim.
Since the trace of a commutator on the invariant space \(W\)
is zero,
\[
 0=\operatorname{tr}_W[x,b]=m\lambda([x,b]).
\]
Characteristic zero gives the asserted vanishing.
If \(w\) is any common \(\lambda\)-eigenvector, then
\(b(xw)=\lambda(b)xw+[b,x]w=\lambda(b)xw\).
This proves stability of the full eigenspace.
\end{proof}

\begin{theorem}[Lie's theorem]\label{res:B02}
Let \(k\) be algebraically closed of characteristic zero.
A nonzero representation of a solvable Lie algebra
\(\mathfrak a\) has a common eigenvector. Every such
representation admits a common upper triangular basis.
In particular, its irreducible representations are
one dimensional.
\end{theorem}

\begin{proof}
Induct on \(\dim\mathfrak a\). For a nonzero solvable
algebra, \([\mathfrak a,\mathfrak a]\ne\mathfrak a\);
a hyperplane \(\mathfrak b\) containing the derived algebra
is an ideal. Write \(\mathfrak a=\mathfrak b+kx\).
Induction gives a common \(\mathfrak b\)-eigenvector and
its character \(\lambda\).
By Lemma~\ref{supp:cyclic-weight}, the common eigenspace
\(V_\lambda\ne0\) is \(x\)-stable. Algebraic closure
supplies an eigenvector of \(x|_{V_\lambda}\).
It is a common eigenvector of \(\mathfrak a\).
Now induct on \(\dim V\), applying the assertion to the
quotient by its eigenline. Lifting the quotient flag
produces a common triangular basis. An irreducible
module must equal its invariant eigenline.
\end{proof}

One consequence needed later is worth proving now. If
\(\mathfrak r\) is a solvable ideal of an arbitrary
\(\mathfrak g\), and \(V\) is an irreducible
\(\mathfrak g\)-module over an algebraically closed
characteristic-zero field, then \(\mathfrak r\) acts by
scalars and \([\mathfrak g,\mathfrak r]\) acts by zero.
Indeed, choose a common \(\mathfrak r\)-eigenvector.
For every \(x\in\mathfrak g\), the cyclic-space lemma for
\(\mathfrak r+kx\) proves
\(\lambda([x,\mathfrak r])=0\) and stability of the same
nonzero eigenspace. It is consequently a
\(\mathfrak g\)-submodule, and therefore is all of \(V\).

\section{Jordan decomposition and the trace criterion}

\begin{proposition}[Jordan decomposition and derivations]
\label{res:B03}
Over a characteristic-zero field every \(A\in\operatorname{End}(V)\)
has a unique decomposition
\[
 A=A_s+A_n,\qquad [A_s,A_n]=0,
\]
where \(A_s\) is semisimple after scalar extension and
\(A_n\) is nilpotent. Both are polynomials in \(A\).
Every \(A\)-stable subspace is stable under both parts.
If \(A=D\) is a derivation of a Lie algebra, then both
\(D_s,D_n\) are derivations.
\end{proposition}

\begin{proof}
Over a splitting field, decompose \(V=\bigoplus V_\lambda\)
into generalized eigenspaces. Choose, by the Chinese
remainder theorem, \(P(t)\) with
\(P(t)\equiv\lambda\pmod{(t-\lambda)^N}\) for every
eigenvalue, where \(N\) is large enough.
Then \(P(A)\) is \(\lambda\) times the identity on
\(V_\lambda\), and \(A-P(A)\) is nilpotent there.
The polynomial can be uniquely chosen modulo the minimal
polynomial of \(A\); its interpolation data are Galois
invariant. In characteristic zero the splitting field is
separable, so this representative has coefficients in \(k\).

For uniqueness, a second commuting decomposition commutes
with \(A\), hence with its polynomial parts. The difference
of the two semisimple parts then equals the opposite
difference of the nilpotent parts. The former is
semisimple (simultaneously diagonalize the commuting
semisimple operators), while the latter is nilpotent
(expand a sufficiently high power of commuting nilpotents).
Both differences are therefore zero. Polynomial expressions
also prove the assertion about invariant subspaces.

Suppose \(D\) is a derivation. For generalized eigenvectors
\(u\in V_\lambda,v\in V_\mu\), the derivation rule gives
\[
 (D-\lambda-\mu)^N[u,v]
 =\sum_{i+j=N}\binom Ni[(D-\lambda)^iu,(D-\mu)^jv].
\]
For large \(N\) every term is zero, so
\([V_\lambda,V_\mu]\subseteq V_{\lambda+\mu}\).
Consequently
\[
 D_s[u,v]=(\lambda+\mu)[u,v]
          =[D_su,v]+[u,D_sv].
\]
By bilinearity this holds everywhere. Thus \(D_s\) is a
derivation, and so is \(D_n=D-D_s\). The identities descend
from the splitting field.
\end{proof}

\begin{theorem}[Cartan's trace criterion]\label{res:B04}
For a Lie subalgebra \(\mathfrak a\subseteq\mathfrak{gl}(V)\)
over a characteristic-zero field, the following are equivalent:
\[
 \mathfrak a\text{ is solvable},\qquad
 \operatorname{tr}(xy)=0
 \quad(x\in[\mathfrak a,\mathfrak a],\ y\in\mathfrak a).
 \tag{2.1}\label{eq:trace-criterion}
\]
In particular, with
\(K_{\mathfrak g}(x,y)=
\operatorname{tr}_{\mathfrak g}(\operatorname{ad}x\operatorname{ad}y)\),
an abstract algebra is solvable if and only if
\(K_{\mathfrak g}([\mathfrak g,\mathfrak g],\mathfrak g)=0\).
\end{theorem}

\begin{proof}
Extension to an algebraic closure preserves traces,
brackets, and vanishing of a term of the derived series.
We may therefore assume \(k\) algebraically closed.
For the forward implication apply Lie's theorem:
derived matrices are strictly upper triangular and the
other matrices upper triangular, so their product has zero
diagonal.

Conversely, fix \(x\in[\mathfrak a,\mathfrak a]\), with
distinct eigenvalues \(\lambda_1,\ldots,\lambda_m\) and
generalized eigenspace dimensions \(d_1,\ldots,d_m\).
Let \(E\) be their finite-dimensional rational span in \(k\)
and choose any rational linear map \(b:E\to\mathbb Q\).
Let \(B\) act by \(b(\lambda_i)\) on the generalized
\(\lambda_i\)-eigenspace of \(x\). We claim
\[
 \operatorname{ad}B=R(\operatorname{ad}x)
 \quad\text{for a polynomial }R\text{ with }R(0)=0.
 \tag{2.2}\label{eq:cartan-polynomial}
\]
Indeed, \(\operatorname{ad}x_s\) has eigenvalues
\(\lambda_i-\lambda_j\) on the corresponding Hom blocks.
Interpolation gives a polynomial \(Q\) with
\(Q(\lambda_i-\lambda_j)=b(\lambda_i)-b(\lambda_j)\)
and \(Q(0)=0\). Thus
\(\operatorname{ad}B=Q(\operatorname{ad}x_s)\).
The identity
\((\operatorname{ad}x)_s=\operatorname{ad}x_s\)
follows by separating the commuting diagonalizable and
nilpotent operators
\(\operatorname{ad}x_s,\operatorname{ad}x_n\).
By Jordan decomposition, the former is \(P(\operatorname{ad}x)\),
where \(P(0)=0\), since zero is an eigenvalue on the identity
endomorphism. Set \(R=Q\circ P\) to obtain
\eqref{eq:cartan-polynomial}.

Write \(x=\sum_j[y_j,z_j]\) in \(\mathfrak a\).
Because \(R\) has no constant term,
\([B,y_j]=R(\operatorname{ad}x)y_j\) belongs to
\([\mathfrak a,\mathfrak a]\). Cyclicity of trace and
\eqref{eq:trace-criterion} give
\[
 \sum_i d_i b(\lambda_i)\lambda_i
 =\operatorname{tr}(Bx)
 =\sum_j\operatorname{tr}([B,y_j]z_j)=0.
\]
Apply the rational linear functional \(b\) to this
identity in \(E\). It yields
\(\sum_i d_i b(\lambda_i)^2=0\) in \(\mathbb Q\).
Each \(d_i\) is a positive integer, so all
\(b(\lambda_i)=0\). As every rational linear functional
has this property, \(E=0\). Hence every eigenvalue of
\(x\) is zero and \(x\) is nilpotent.

Engel's theorem now makes \([\mathfrak a,\mathfrak a]\)
nilpotent, hence solvable. Its abelian quotient proves
solvability of \(\mathfrak a\).
For the abstract assertion apply the matrix criterion
to the adjoint image \(\mathfrak g/Z(\mathfrak g)\);
the missing kernel is central and therefore solvable.
\end{proof}

\section{The radical and semisimple structure}

\begin{theorem}[The radical and the Killing criterion]\label{res:B05}
The sum of all solvable ideals is a solvable ideal
\(\operatorname{rad}\mathfrak g\), and
\(\mathfrak g/\operatorname{rad}\mathfrak g\) has no nonzero
solvable ideal. In characteristic zero,
\(\operatorname{rad}\mathfrak g=0\) if and only if the Killing
form is nondegenerate. Such an algebra is called semisimple.
\end{theorem}

\begin{proof}
For solvable ideals \(I,J\), the quotient
\((I+J)/I\simeq J/(I\cap J)\) is solvable, so \(I+J\)
is solvable. A finite-dimensional sum is already a finite
sum, proving the first assertion. A solvable ideal in the
quotient lifts to a solvable extension of the radical,
so must be zero.

Trace cyclicity gives symmetry and invariance of the Killing
form:
\[
 K([x,y],z)=K(x,[y,z]).
 \tag{2.3}\label{eq:killing-invariant}
\]
Its kernel \(I\) is an ideal. For any ideal \(J\) and
\(x,y\in J\), the action of \(x,y\) on \(\mathfrak g/J\)
is zero, so a block triangular trace calculation gives
\[
 K_J(x,y)=K_{\mathfrak g}(x,y).
 \tag{2.4}\label{eq:killing-ideal}
\]
If \(\mathfrak g\) is semisimple, \(K_I=0\).
The trace criterion makes \(I\) solvable; hence \(I=0\).

Conversely, suppose \(K_{\mathfrak g}\) is nondegenerate
and the radical is nonzero. Its last nonzero derived term
is a nonzero abelian ideal \(A\) of \(\mathfrak g\).
For \(a\in A,y\in\mathfrak g\),
\(\operatorname{ad}a\operatorname{ad}y\) sends
\(\mathfrak g\) into \(A\) and kills \(A\):
\(\operatorname{ad}y\) preserves \(A\), and
\([a,A]=0\). Its square is zero, so its trace is zero.
Thus \(A\) lies in the Killing kernel, a contradiction.
\end{proof}

This also proves that scalar extension preserves
semisimplicity: its Killing matrix is the same matrix
over the larger field and remains invertible. We shall
use this consequence when passing from real algebras
to their complexifications.

\begin{theorem}[Simple-ideal decomposition]\label{res:B06}
A finite-dimensional characteristic-zero semisimple algebra
is a direct sum of nonabelian simple ideals. Every ideal
has an ideal complement, namely its Killing orthogonal
complement. Its center is zero and
\([\mathfrak g,\mathfrak g]=\mathfrak g\).
Every ideal is a sum of a subset of its simple ideals.
\end{theorem}

\begin{proof}
If \(I\) is an ideal, invariance shows that \(I^\perp\)
is an ideal. Set \(J=I\cap I^\perp\). By
\eqref{eq:killing-ideal}, its Killing form is zero,
since \(J\subseteq I^\perp\) and \(J\subseteq I\).
Thus \(J\) is solvable, and semisimplicity forces \(J=0\).
Nondegeneracy and dimensions then give
\(\mathfrak g=I\oplus I^\perp\).
The two ideals commute because their bracket lies in
their zero intersection.

Take a minimal nonzero ideal and split it off.
An ideal of this summand is an ideal of the full direct
sum, so minimality makes the summand simple; it cannot
be abelian. Iterating decreases dimension and gives the
decomposition. A simple nonabelian algebra is perfect,
since its derived algebra is a nonzero ideal, and its
center is zero. These assertions pass to direct sums.

Write \(\mathfrak g=\bigoplus_i\mathfrak g_i\) and let
\(I\) be any ideal. Its projection to \(\mathfrak g_i\)
is either zero or all of that simple algebra. In the
latter case
\([\mathfrak g_i,I]=
[\mathfrak g_i,\operatorname{pr}_i I]=\mathfrak g_i\),
so \(\mathfrak g_i\subseteq I\). Removing these
summands leaves zero. This proves the final assertion
and uniqueness of the complementary sum.
\end{proof}

\begin{theorem}[Inner derivations and abstract Jordan decomposition]
\label{res:B07}
Every derivation of a characteristic-zero semisimple
\(\mathfrak g\) is inner. Each \(x\in\mathfrak g\) has a
unique expression
\[
 x=x_s+x_n,\qquad [x_s,x_n]=0,
\]
where \(\operatorname{ad}x_s\) is semisimple and
\(\operatorname{ad}x_n\) is nilpotent. If \([x,y]=0\),
then \([x_s,y]=[x_n,y]=0\).
\end{theorem}

\begin{proof}
Inside \(\operatorname{Der}\mathfrak g\), the adjoint
image is an ideal, because
\([D,\operatorname{ad}x]=\operatorname{ad}(Dx)\).
On derivations use the invariant trace form
\(B(D,E)=\operatorname{tr}_{\mathfrak g}(DE)\).
Its restriction to \(\operatorname{ad}\mathfrak g\)
is the nondegenerate Killing form. Hence
\[
 \operatorname{Der}\mathfrak g
 =\operatorname{ad}\mathfrak g\oplus
   (\operatorname{ad}\mathfrak g)^\perp.
\]
The second summand is an ideal by invariance.
It commutes with the first, since their intersection
is zero. For \(D\) in this complement,
\(\operatorname{ad}(Dx)=[D,\operatorname{ad}x]=0\).
The center is zero, so \(Dx=0\) for all \(x\), and
the complement vanishes.

Apply Proposition~\ref{res:B03} to
\(\operatorname{ad}x\) on \(\mathfrak g\).
Both matrix Jordan parts are derivations, so are
\(\operatorname{ad}x_s,\operatorname{ad}x_n\)
for uniquely determined elements. Their sum and
commutator identify \(x=x_s+x_n\) and
\([x_s,x_n]=0\), using injectivity of the adjoint
map. Matrix uniqueness gives uniqueness here.
If \([x,y]=0\), then \(\operatorname{ad}y\)
commutes with \(\operatorname{ad}x\) and its
polynomial parts. Hence the two asserted brackets
are central, and therefore zero.
\end{proof}

\begin{theorem}[Characteristic radical and nilradical]\label{res:B08}
In characteristic zero the radical and the largest nilpotent
ideal \(\operatorname{nilrad}\mathfrak g\) are preserved by
every derivation. More precisely,
\[
 D(\operatorname{rad}\mathfrak g)
       \subseteq\operatorname{nilrad}\mathfrak g
 \quad(D\in\operatorname{Der}\mathfrak g),\qquad
 [\mathfrak g,\operatorname{rad}\mathfrak g]
       \subseteq\operatorname{nilrad}\mathfrak g.
 \tag{2.5}\label{eq:nilrad-characteristic}
\]
The nilradical consists exactly of those \(x\) in the
radical for which \(\operatorname{ad}_{\mathfrak g}x\)
is nilpotent.
\end{theorem}

\begin{proof}
We first establish the useful equality
\[
 \operatorname{rad}\mathfrak g
     =[\mathfrak g,\mathfrak g]^\perp
 \tag{2.6}\label{eq:rad-perp}
\]
under the Killing form. Over an algebraic closure,
choose a composition series of the adjoint module.
The consequence of Lie's theorem proved after
Theorem~\ref{res:B02} says that the radical acts
by a scalar on each simple quotient. A commutator
has zero trace on each quotient. Thus
\(K(\operatorname{rad}\mathfrak g,
[\mathfrak g,\mathfrak g])=0\).
For the reverse containment let
\(J=[\mathfrak g,\mathfrak g]^\perp\).
It is an ideal by invariance. If \(a,b,c\in J\),
then \([a,b]\in[\mathfrak g,\mathfrak g]\), so
\(K([a,b],c)=0\). By \eqref{eq:killing-ideal},
the trace criterion makes \(J\) solvable.
It therefore lies in the radical.
The same argument proves the equality over \(k\):
one inclusion is checked after extension, and for
the other \(J\) is already defined over \(k\).
In particular this equality shows that the radical
commutes with scalar extension.

For a derivation \(D\), the identity
\([D,\operatorname{ad}x]=\operatorname{ad}(Dx)\)
and trace cyclicity give
\[
 K(Dx,y)+K(x,Dy)=
 \operatorname{tr}[D,\operatorname{ad}x\operatorname{ad}y]=0.
\]
Since \(D\) preserves the derived algebra,
\eqref{eq:rad-perp} proves \(D\mathfrak r\subseteq\mathfrak r\),
where \(\mathfrak r=\operatorname{rad}\mathfrak g\).

Over an algebraic closure, triangularize the action of
\(\mathfrak r\) on \(\mathfrak g\). Let \(\mathfrak n\)
be the subspace on which all its diagonal characters
vanish. Equivalently,
\[
 \mathfrak n=\{x\in\mathfrak r:
             \operatorname{ad}_{\mathfrak g}x
             \text{ is nilpotent}\}.
\]
For any derivation \(D\), the semidirect algebra
\(kD\ltimes\mathfrak r\) is solvable and acts on
\(\mathfrak g\) through \(D\) and
\(\operatorname{ad}\mathfrak r\). Lie's theorem
triangularizes this whole action. Its commutators,
in particular
\([D,\operatorname{ad}x]=\operatorname{ad}(Dx)\)
for \(x\in\mathfrak r\), are strictly upper
triangular. Thus \(D\mathfrak r\subseteq\mathfrak n\).
Taking \(D=\operatorname{ad}y\) shows that
\(\mathfrak n\) is an ideal. Engel applied to its
adjoint action on itself makes it nilpotent.

If \(A\) is any nilpotent ideal, it is solvable,
so lies in \(\mathfrak r\). For \(a\in A\),
\(\operatorname{ad}a\) sends \(\mathfrak g\) into
\(A\), and its further powers send \(A\) down
the lower central series of \(A\).
It is nilpotent on \(\mathfrak g\), so \(A\subseteq
\mathfrak n\). This proves maximality.
Finally \(D\mathfrak n\subseteq D\mathfrak r
\subseteq\mathfrak n\).

Over a nonclosed field, define \(\mathfrak n\) by
the same nilpotent-adjoint condition. The preceding
triangular calculation after scalar extension proves
that its \(k\)-rational elements form a subspace and
an ideal, and proves all the stated derivation
inclusions there. Engel and the maximality argument
are already valid over \(k\). Hence the conclusion
descends without choosing eigenvalues in \(k\).
\end{proof}

The trace criterion, rather than a representation
classification, supplied every semisimple fact used
in this chapter. The complete reducibility of arbitrary
finite-dimensional representations is a separate
theorem, proved next.
